%AttackName: M-GAP UAP-1

%Input:
Let \( f \) be the target model, \( x \) be the original input data point, \( y \) be the true label, and \( \epsilon \) be the maximum allowable perturbation. The GAP-UAP Attack aims to generate a universal perturbation \( \delta \) that effectively deceives the model by leveraging the gradient information.

%Output:
The output of the GAP-UAP Attack is a universal adversarial perturbation \( \delta \) that can be applied to various inputs to mislead the classifier.

%Formula:
The GAP-UAP Attack can be formulated as follows:
1. Initialize the universal perturbation \( \delta \):
   $
   \delta = 0.
   $
2. For each iteration \( t \):
   $
   \delta^{(t)} = \delta^{(t-1)} + \alpha \cdot \text{sign} \left( \frac{1}{N} \sum_{i=1}^N \nabla_{\delta} L(f(x_i + \delta^{(t-1)}), y_i) \right),
   $
   where \( N \) is the number of sampled inputs and \( L \) is the loss function designed to maximize the misclassification probability.
3. Project \( \delta \) to ensure it remains within the allowable perturbation bounds:
   $
   \delta = \text{clip}(\delta, -\epsilon, \epsilon).
   $

%Explanation:
The Gradient-Averaged Perturbation Universal Adversarial Perturbation (GAP-UAP) Attack generates a universal adversarial perturbation \( \delta \) by utilizing gradient information from a set of inputs. By optimizing \( \delta \) based on the average gradient of the loss function over multiple samples, the attack creates a perturbation that is effective across a range of inputs. This method illustrates the potential for gradient-based approaches to craft universal adversarial perturbations that exploit the vulnerabilities of machine learning models, demonstrating the efficacy of such attacks in practical scenarios.